\subsection{Altre tecnologie}
\label{subsec:altre_tecnologie}

\subsubsection{Docker}
\label{subsubsec:docker}
Docker è un software open source che consente lo sviluppo e la distribuzione di applicazioni all'interno di container, ambienti isolati che condividono il kernel del sistema operativo.\\  

\textbf{Motivazioni:}  
\begin{enumerate}
    \item \textbf{Isolamento}: ogni applicazione viene eseguita all'interno di un container dedicato, che rappresenta un ambiente di esecuzione isolato.  
    Questo evita problemi di incompatibilità tra le dipendenze delle diverse applicazioni, comprese quelle in esecuzione sulla macchina target.  
    Inoltre, eventuali malfunzionamenti rimangono confinati all'interno del container, senza compromettere l'intero sistema;
    \item \textbf{Portabilità}: l'ambiente di sviluppo è identico a quello di produzione, questo rende la transizione semplice e senza errori inaspettati.
    Infatti basta che sulla macchina target sia installato Docker; 
    \item \textbf{Semplificazione dello sviluppo}: consente di automatizzare la configurazione e l'esecuzione dell'ambiente di runtime attraverso un insieme di comandi strutturato.  
    Ciò riduce la probabilità di errori e semplifica l'avvio del sistema a un solo comando;
    \item \textbf{Scalabilità}: i container sono più leggeri rispetto alle macchine virtuali, questo permette di aggiungere e rimuovere container in modo rapido ed efficiente.
\end{enumerate}

\subsubsection{PostgreSQL}
\label{subsubsec:postgresql}
PostgreSQL è un database relazionale open source avanzato, noto per la sua affidabilità, scalabilità e conformità agli standard SQL.\\

\textbf{Motivazioni:}
\begin{enumerate}
    \item \textbf{Open Source}: PostgreSQL è completamente open source, senza costi di licenza, con una comunità attiva che lo mantiene e aggiorna continuamente.
    \item \textbf{Affidabilità e integrità dei dati}: Supporta transazioni ACID (Atomicità, Consistenza, Isolamento, Durabilità), garantendo la sicurezza e la coerenza dei dati.
\end{enumerate}

\subsubsection{Hugging Face}
\label{subsubsec:huggingface}
Hugging Face è una piattaforma leader nell'intelligenza artificiale che fornisce strumenti open-source per il Natural Language Processing (NLP), la visione artificiale e altri campi dell'AI. 
Offre una vasta libreria di modelli pre-addestrati, dataset ottimizzati e API facili da usare, supportando framework come PyTorch.\\

\textbf{Motivazioni:}
\begin{enumerate}
    \item \textbf{Comunità attiva e collaborativa}: Una vasta community di sviluppatori e ricercatori contribuisce con modelli, dataset e miglioramenti costanti.
    \item \textbf{Compatibilità con diverse piattaforme }: Supporta framework come TensorFlow, PyTorch e JAX, rendendo l'adozione flessibile per diversi ambienti di sviluppo.
    \item \textbf{Facilità d'uso}: Le API e le librerie semplificano l'integrazione di modelli avanzati con poche righe di codice.
    \item \textbf{Ampia libreria di modelli pre-addestrati}: Offre migliaia di modelli di intelligenza artificiale per NLP, visione artificiale e altre applicazioni, riducendo i tempi di sviluppo.
    \item \textbf{Strumenti di valutazione e benchmarking}: Fornisce metriche integrate per confrontare le prestazioni dei modelli e scegliere le migliori soluzioni.
\end{enumerate}